\textbf{Datum}: 17/12/2024 \newline
\textbf{Chair}: Alex Guo \newline
\textbf{Secretary}: Gerrald van Weelden \newline

\subsubsection{\textbf{Agenda}}
\begin{enumerate}
    \item Opening 
    \item Aannemen notulen 
    \item Ingekomen stukken
    \item Deadlines 
    \item Update per groep 
    \item WVTTK 
    \item Rondvraag 
    \item Afsluiting 
\end{enumerate}
\subsubsection{\textbf{Persons present}}
\begin{itemize}
    \item Alex Guo 
    \item Gerrald van Weelden
    \item Itai Givony 
    \item Karsten Smid 
    \item Luuk Melissant 
    \item Niek van Vroonhoven 
    \item Noah Herijgers 
    \item Stijn ter Huurne 
    \item Teun Bello
\end{itemize}
\subsubsection{\textbf{Minutes}}
\begin{enumerate}
    \item Opening \newline
    \begin{itemize}
        \item Alex opent de vergadering.
    \end{itemize}
    \item Aannemen notulen 
    \begin{itemize}
        \item Alex vraagt naar eventuele bezwaren tegen notulen van 05/12/2024. Er zijn geen bezwaren.
    \end{itemize}
    \item Ingekomen stukken 
    \begin{itemize}
        \item Gerrald geeft aan dat er geen ingekomen stukken zijn.
    \end{itemize}
    \item Deadlines 
    \begin{itemize}
        \item Alex noemt de belangrijkste deadlines deze week: buddycheck woensdag 22:00 uur en Intermediate Report op vrijdag.
        \item Stijn vraagt welke tijd de report ingeleverd moet zijn. Alex geeft aan dat het vrijdag 24.00 uur middernacht is.
    \end{itemize}
    \item Update per groep 
    \begin{itemize}
        \item Niek stelt voor dat de filter samen presenteren, wordt aangenomen
        \item Alex heeft niks extra’s meer te vertellen over de Power Supply, gaan deze week verder werken aan Intermediate Report.
        \item Stijn vraagt naar de conclusie, Alex antwoordt dat het wel gaat lukken, hoewel Gerrald het deze week wel druk heeft.
        \item Itai vraagt om een conclusie wat we met de Booming Bass gaan doen. Daarop antwoordt Gerrald dat hij nog bezig is met het uitrekenen van de componenten en dat het waarschijnlijk een Linkwitz transform wordt. Itai vraagt om verduidelijking of het doorgaat, Gerrald antwoordt daarop dat hij voornemens is het te doen.
        \item Power Amplifier presenteert. Vorige keer hadden ze een te hoge cut-off frequency van 47 kHz. 2 uur lang vorige week getest, kwamen er achter dat het waarschijnlijk aan de condensatoren ligt. De cut-off frequency is nu 41.8 kHz. De Power Amplifier is dus in principe klaar, deze week gaan ze aan het verslag werken en metingen doen.
        \item Stijn en Hendrik vragen opheldering waarom ze weten dat de tweede meting wel goed is. Daarop antwoordt Itai dat de condensatoren eerst in female headers zaten, maar nu vast gesoldeerd zijn en dat het daarom nu wel werkt. Hendrik vraagt of het meerdere keren getest is. Itai bevestigd dat het vorige week meerdere keren getest is, een professor heeft geholpen met bode plots op de oscilloscoop vorige week en dat ook hij tot de conclusie kwam dat het aan de condensatoren lag.
        \item Teun vraagt of het probleem van vorige keer met een amplification van 2.5 i.p.v. 25 opgelost is. Daarop antwoordt Itai dat het lag aan een instelling van de oscilloscoop, wat vorige week snel opgelost was. 
        \item Stijn vraagt hoe het gaat met de report. Itai en Luuk geven aan het wel te redden, hoewel het nog best wat werk is. Stijn moedigt aan vroegtijdig een seintje te geven als het niet gaat lukken.
        \item Filters presenteren. Ze hadden het probleem dat de electrical respons vorige week niet plat was. Deze week is het nog getest, en de akoestische respons is vrijwel plat, alleen de tweeter is wat hoger, kan waarschijnlijk verholpen worden met een weerstandje. 
        \item Karsten merkt op dat de microfoon gemikt stond op de tweeter, dus het kan zijn dat de piek daardoor niet te zien was. Stijn merkt op dat de speaker ook getest kan worden met de microfoon verder van de speaker af, waardoor het probleem kleiner zal zijn. Teun merkt ook op dat het met het oor tijdens de meting ook een witte ruis leek te zijn en het dus misschien ook niet te horen is tijdens het afspelen van muziek.
        \item Intermediate report wordt besproken. Er wordt afgesproken het vanavond af proberen te hebben om nog feedback te krijgen van Hendrik. Hendrik geeft aan dat hij in deze korte tijd niet veel feedback kan geven en de feedback waarschijnlijk globaal zal zijn. Stijn vraagt of Hendrik aandacht kan besteden aan de grafieken of de vorm goed is. Hendrik antwoord bevestigend. 
        \item Teun sluit af voor de filters.
        \item Karsten vraagt of de data voor de filters over de design in het report opgeschreven kan worden, met formules o.i.d. Teun geeft aan de Python code nog op te zoeken. Stijn zegt zijn spullen ook wel op te zoeken. 
    \end{itemize}

    \item WVTTK 
    \begin{itemize}
        \item Gerrald geeft aan dat er niets ter tafel gekomen is.
    \end{itemize}
    \item Rondvraag 
    \begin{itemize}
        \item Alex vraagt of er nog mensen zijn die iets willen inbrengen in de vergadering. Dat blijkt niet zo te zijn.
    \end{itemize}
    \item Afsluiting 
    \begin{itemize}
        \item Alex geeft aan dat dat het was en besluit daarmee de vergadering.
    \end{itemize}
\end{enumerate}